\section{Arbitrarily small nonhomogeneous perturbations can remove the gap}
\label{sec:r9-bounded-boundary}

The upper dilation condition has a substantive role even among symbols
of the same order and arbitrarily close in quadratic-form norm. The
following construction complements the exclusion theorem. Its existence
proof uses standard constrained minimization; no separate priority claim
is attached to that variational method.

Fix a principal symbol $a$ as in the cone theorem, of degree $m>1$,
and write
\[
 r=m/N,\qquad p=2+2r,\qquad M_c=M_{c,a},\qquad
 A(f)=\int a(\xi)|\widehat f(\xi)|^2\,d\xi.
\]
For fixed numbers $\varepsilon,q>0$, set
\begin{equation}\label{eq:r9-bounded-symbol}
 b_{\varepsilon,q}(\xi)
 =\varepsilon a(\xi)\frac{|\xi|^q}{1+|\xi|^q},\qquad
 B(f)=\int b_{\varepsilon,q}(\xi)|\widehat f(\xi)|^2\,d\xi,
 \qquad E(f)=\frac12[A(f)+B(f)]-\frac1p\|f\|_p^p.
\end{equation}
The symbol is real, even, individually $C^{N+2}$ away from zero, and
satisfies all the order-$m$ derivative bounds of the cone theorem. In
particular,
\begin{equation}\label{eq:r9-bounded-formsmall}
 0\le b_{\varepsilon,q}\le\varepsilon a,
 \qquad 0\le B(f)\le\varepsilon A(f).
\end{equation}
But its dilation defect has the opposite sign:
\begin{equation}\label{eq:r9-bounded-defect}
 mb_{\varepsilon,q}-\xi\cdot\nabla b_{\varepsilon,q}
 =-\varepsilon q\,a(\xi)
       \frac{|\xi|^q}{(1+|\xi|^q)^2}<0\quad(\xi\ne0).
\end{equation}

The number $M_c$ is also the actual sharp GN threshold for the full
kinetic form $A+B$, not merely a comparison with the principal part.
Indeed $A+B\ge A$ gives the principal sharp constant as an upper
bound for the full one. If $Q$ is a principal optimizer and
$Q_\lambda(x)=\lambda^{N/2}Q(\lambda x)$, dominated convergence gives
\[
 \frac{B(Q_\lambda)}{\lambda^m}
 =\varepsilon\int a(\xi)
 \frac{\lambda^q|\xi|^q}{1+\lambda^q|\xi|^q}
 |\widehat Q(\xi)|^2\,d\xi\longrightarrow0.
\]
The sharp constants therefore coincide. The full inequality has no
nonzero optimizer: $B(f)>0$ makes its comparison with the principal
sharp inequality strict for every nonzero $f$.

\begin{theorem}[A boundary within the same operator order]
\label{thm:r9-bounded-boundary}
Define
\[
 I(M)=\inf\{E(f): f\in H^{m/2}(\mathbb R^N;\mathbb C),\ \|f\|_2^2=M\}.
\]
For $0<M\le M_c$, $I(M)=0$ and the infimum is not attained.
For every
\begin{equation}\label{eq:r9-bounded-window}
 M_c<M<(1+\varepsilon)^{N/m}M_c,
\end{equation}
$-\infty<I(M)<0$, and each minimizing sequence has, after translations
and passage to a subsequence, a strong $H^{m/2}$ limit attaining the
infimum. A minimizer $Q_M$ can be chosen real and belongs to
$L^\infty\cap H^{m+1}$. For some $\omega_M>0$ it solves
\begin{equation}\label{eq:r9-bounded-profile}
 [a(D)+b_{\varepsilon,q}(D)]Q_M+\omega_M Q_M
       -|Q_M|^{2r}Q_M=0.
\end{equation}
Its energy is exactly
\begin{equation}\label{eq:r9-bounded-energy}
 E(Q_M)=-\frac{\varepsilon q}{2m}
 \int a(\xi)\frac{|\xi|^q}{(1+|\xi|^q)^2}
             |\widehat Q_M(\xi)|^2\,d\xi<0.
\end{equation}
Thus the modified evolution has a real-profile, zero-momentum standing
wave $u_M(t)=e^{i\omega_Mt}Q_M$ for every mass in
\eqref{eq:r9-bounded-window}. For each fixed $\varepsilon>0$, these
negative-energy standing-wave masses approach $M_c$ from above.
\end{theorem}

\begin{proof}
The sharp principal inequality is
\begin{equation}\label{eq:r9-bounded-GN}
 \frac1p\|f\|_p^p\le
 \frac12(M(f)/M_c)^r A(f).
\end{equation}
Consequently, if $0<M\le M_c$ then $E(f)\ge B(f)/2>0$ for every
nonzero admissible $f$. The strict inequality holds because
$b_{\varepsilon,q}>0$ away from the single frequency zero. A smooth
mass-$M$ trial function dilated to low frequency has energy tending to
zero, so $I(M)=0$ without attainment.

Now fix $M$ in \eqref{eq:r9-bounded-window}, and put
$d=(M/M_c)^r-1\in(0,\varepsilon)$. For every
$\eta\in(0,\varepsilon)$ there is a finite $C_\eta$ such that
\[
 b_{\varepsilon,q}(\xi)\ge(\varepsilon-\eta)a(\xi)-C_\eta.
\]
Indeed the coefficient $\varepsilon|\xi|^q/(1+|\xi|^q)$ tends to
$\varepsilon$ at infinity, and $a$ is bounded on every finite ball.
Choose $\eta<\varepsilon-d$. Equation~\eqref{eq:r9-bounded-GN} gives
\begin{equation}\label{eq:r9-bounded-coercivity}
 E(f)\ge\frac{\varepsilon-\eta-d}{2}A(f)-\frac{C_\eta M}{2}
 \quad\text{when }M(f)=M.
\end{equation}
Thus $I(M)>-\infty$, and minimizing sequences are bounded in
$H^{m/2}$.

To see strict negativity, multiply a real principal GN optimizer of
mass $M_c$ by $\sqrt{M/M_c}$ and approximate it in $H^{m/2}$ by real
$C_c^\infty$ functions, adjusting the mass. This yields a real smooth
$f$ of mass exactly $M$ with
$E_a(f)=A(f)/2-\|f\|_p^p/p<0$.
For $f_\lambda(x)=\lambda^{N/2}f(\lambda x)$,
\[
 B(f_\lambda)\le\varepsilon\lambda^{m+q}
       \int a(\xi)|\xi|^q|\widehat f(\xi)|^2\,d\xi,
 \qquad E(f_\lambda)=\lambda^m E_a(f)+\tfrac12B(f_\lambda)<0
\]
for sufficiently small positive $\lambda$. This proves $I(M)<0$.
The same statements hold at every smaller positive mass in the window
where needed below.

We give the compactness argument to handle the mass-critical power
without presuming that the added symbol has higher order.
For $\theta>1$ such that the masses under consideration remain below
the upper endpoint of \eqref{eq:r9-bounded-window}, amplitude scaling
gives
\[
 I(\theta L)\le\theta I(L),\qquad
 I(\theta L)\le\theta^{p/2}I(L)<\theta I(L)
       \quad\text{if }I(L)<0.
\]
The second inequality follows from
\[
 E(\sqrt\theta f)=\theta^{p/2}E(f)
       +\tfrac12(\theta-\theta^{p/2})[A(f)+B(f)].
\]
For every split $M=M_1+M_2$ with $M_j>0$, this implies
\begin{equation}\label{eq:r9-bounded-binding}
 I(M)<I(M_1)+I(M_2).
\end{equation}
If both pieces have nonnegative infimum, their infima are zero while
$I(M)<0$. Otherwise $I(M)/M$ is strictly below the energy-per-mass
ratio of a negative piece and no greater than that of the other piece;
weighting by $M_1,M_2$ proves \eqref{eq:r9-bounded-binding}.

A bounded minimizing sequence $f_n$ cannot vanish locally in mass.
Bounded-overlap localization of the GN inequality gives
\[
 \|f_n\|_p^p
 \le C\left(\sup_y\int_{B(y,2)}|f_n|^2\right)^r
           \|f_n\|_{H^{m/2}}^2.
\]
For example, apply the GN inequality to smooth unit-cube cutoffs and
sum their squared $H^{m/2}$ norms, which are bounded by a constant
times the global squared norm. Local vanishing would imply
$\|f_n\|_p\to0$, contradicting $I(M)<0$.
After translations and extraction, there is therefore a nonzero weak
limit $Q$ in $H^{m/2}$, with almost-everywhere convergence. Put
$w_n=f_n-Q$. The two quadratic forms split by weak Hilbert convergence
(the $b$ form is bounded by $\varepsilon A$), and the nonlinear
integral splits by the Brezis--Lieb lemma. If $0<M(Q)<M$, normalizing
$w_n$ to mass exactly $M-M(Q)$ changes its energy by $o(1)$, because
its energy terms are bounded. The splitting would then contradict
\eqref{eq:r9-bounded-binding}. Hence $M(Q)=M$ and $M(w_n)\to0$.
The principal GN inequality gives
\[
 \|w_n\|_p^p\le C M(w_n)^r A(w_n)\longrightarrow0.
\]
Energy splitting and $E(Q)\ge I(M)$ now force
$A(w_n)+B(w_n)\to0$. This proves strong $H^{m/2}$ convergence and
attainment.

A real minimizer is available without any modulus inequality for
higher-order kinetic forms. Write $f=X+iY$ and, when both real
components have positive mass, set
$\vartheta=M(X)/M(f)$, $g=X/\sqrt\vartheta$ and
$h=Y/\sqrt{1-\vartheta}$. Both $g,h$ have mass $M(f)$. Real-evenness
splits the kinetic form, whereas convexity gives
\[
 |f|^p=[\vartheta g^2+(1-\vartheta)h^2]^{p/2}
       \le\vartheta|g|^p+(1-\vartheta)|h|^p.
\]
Thus $E(f)\ge\vartheta E(g)+(1-\vartheta)E(h)$, and the complex
and real constrained infima coincide. A real minimizing sequence
and the proved compactness yield a real minimizer.

The Lagrange multiplier rule gives \eqref{eq:r9-bounded-profile}.
Testing it against $Q_M$ shows
\[
 \omega_M M=\|Q_M\|_p^p-A(Q_M)-B(Q_M)
       =\left(1-\frac2p\right)\|Q_M\|_p^p-2E(Q_M)>0.
\]
The symbol $a+b_{\varepsilon,q}+\omega_M$ is comparable with
$1+|\xi|^m$ and its resolvent satisfies the individual Mihlin bounds.
The finite Sobolev bootstrap for the principal critical power therefore
gives $L^\infty$, then $H^m\subset H^1$, and finally $H^{m+1}$ by
the local $C^1$ Sobolev chain rule. This is the same individual
regularity argument as for the cone theorem and uses no sign of the
dilation defect.

Finally the mass-preserving scaling curve through the minimizer has
its scalar energy minimum at $\lambda=1$. Differentiation of its
explicit Fourier energy formula is justified by
$|\xi\cdot\nabla b_{\varepsilon,q}|\le(m+q)b_{\varepsilon,q}$
and domination by $a$. It yields
\[
 \frac m2 A(Q_M)
 +\frac12\int\xi\cdot\nabla b_{\varepsilon,q}(\xi)
                         |\widehat Q_M(\xi)|^2\,d\xi
 -\frac mp\|Q_M\|_p^p=0.
\]
Subtracting this identity divided by $m$ from the energy formula
proves \eqref{eq:r9-bounded-energy}. No spatial moment or
differentiability of the scaling curve in the energy norm is needed.
A real $Q_M$ has zero momentum, and its standing-wave formula supplies
a global solution directly.
\end{proof}

For the original fourth-order principal symbol, one concrete example is
\[
 a(\xi)=|\xi|^4,\qquad
 b_{\varepsilon,2}(\xi)=\varepsilon\frac{|\xi|^6}{1+|\xi|^2}.
\]
Its perturbation is at most $\varepsilon$ times the original kinetic
form, yet its negative-energy zero-momentum standing waves have every
mass in $(M_c,(1+\varepsilon)^{N/4}M_c)$. This is a result about the
specified modified equation, not a counterexample for the pure symbol
$|\xi|^4$.

The construction does not say that every symbol violating the upper
dilation inequality loses its gap. It proves that removing that
condition admits counterexamples arbitrarily close to the principal
operator, even with the same energy space and the same order bounds.
In particular, the homogeneous-symbol robustness theorem cannot be
extended to all small nonhomogeneous quadratic-form perturbations.
The distinction is structural: smallness alone does not control the
sign of the dilation defect.
